% <<vi:Thư xin việc / thư trang trọng — tạo bằng latexkit. Biên dịch: latexmk (kết quả ở build/main.pdf).||en:Cover / formal letter — made with latexkit. Build: latexmk (output in build/main.pdf).>>
% <<vi:Thay các đoạn trong ngoặc vuông bằng nội dung thật. Giữ thư trong 1 trang.||en:Replace the bracketed paragraphs with real text. Keep the letter to 1 page.>>
\documentclass[a4paper,11pt]{article}
\usepackage[a4paper, left=2.3cm, right=2.3cm, top=1.5cm, bottom=2cm]{geometry}
\usepackage[{{PKGOPTS}}]{latexkit}
\usepackage{latexkit-cv}
\pagestyle{empty}
\setlength{\parindent}{0pt}
\setlength{\parskip}{0.7em}

\begin{document}

\letterheader{{{AUTHOR}}}
  {\href{mailto:email@example.com}{email@example.com},
   \href{https://linkedin.com/in/username}{linkedin.com/in/username},
   +84 900 000 000, <<vi:TP. Hồ Chí Minh||en:City, Country>>}

\hfill <<vi:TP. Hồ Chí Minh, ngày \today||en:\today>>

<<vi:Kính gửi: Ban Tuyển dụng Công ty XYZ||en:Hiring Team\\ XYZ Company\\ 123 Example Street, City>>

\textbf{<<vi:V/v: Ứng tuyển vị trí [tên vị trí]||en:Re: Application for [position title]>>}

%%IF en
Dear Hiring Manager,
%%ENDIF

<<vi:[Đoạn 1: vị trí ứng tuyển, nơi biết đến tin tuyển dụng và lý do quan tâm đến công ty.]||en:[Paragraph 1: the position, where the posting was found and why the company is of interest.]>>

<<vi:[Đoạn 2: hai hoặc ba kinh nghiệm, dự án phù hợp nhất với yêu cầu công việc, kèm kết quả cụ thể.]||en:[Paragraph 2: two or three experiences or projects that best match the requirements, with concrete results.]>>

<<vi:[Đoạn 3: điều có thể đóng góp cho nhóm và lý do phù hợp với văn hóa công ty.]||en:[Paragraph 3: what can be contributed to the team and why it fits the company culture.]>>

<<vi:[Đoạn 4: cảm ơn, mong có buổi phỏng vấn, nhắc CV đính kèm.]||en:[Paragraph 4: thanks, hope for an interview, mention the attached CV.]>>

\vspace{1em}
\begin{flushright}
  <<vi:Trân trọng,||en:Sincerely,>>\\[2.5em]
  {{AUTHOR}}
\end{flushright}

\end{document}
